%! AP Statistics — Unit 7, Lesson 7.1 — Group Activity
\documentclass[10pt]{article}
\usepackage{apstats-article}
\usepackage{apstats-boxes}

\begin{document}

\pageheader{Unit 7, Lesson 7.1}{Group Activity}
\namepartnerperiod

\vspace{0.1in}
\textbf{Directions:} Work with your partner on each scenario. For each, (a) state the
hypotheses, (b) describe a Type I error in context, (c) describe a Type II error in
context, and (d) decide which error is more costly and explain why.

\begin{tcolorbox}[colback=white, colframe=black!40, title=\textbf{Tier R --- Remediate},
    fonttitle=\bfseries, arc=2mm, left=3mm, right=3mm, top=2mm, bottom=2mm]

\textbf{Scenario 1 --- School Nurse Study}\\
A school nurse believes the mean number of sick days per year for students at her school
differs from the district average of 5 days.

\begin{enumerate}[label=\alph*.]
  \item State the hypotheses in words (they are given as a starting point):
        $H_0$: The mean number of sick days is \blank{3cm} days.
        $H_a$: The mean number of sick days is \blank{4cm} 5 days.

        \vspace{0.05in}

  \item Describe a Type I error in context (use the sentence frame):
        ``Concluding the mean number of sick days \blank{5cm} when it
        actually \blank{5cm}.''
        \writelines{1}

  \item Describe a Type II error in context:
        ``Failing to conclude \blank{5cm} when it actually \blank{5cm}.''
        \writelines{1}

  \item Which error is more costly? Circle one: \quad \textbf{Type I} \quad \textbf{Type II}
        \writelines{1}
\end{enumerate}
\end{tcolorbox}

\begin{tcolorbox}[colback=white, colframe=black!40, title=\textbf{Tier A --- Approaching Proficiency},
    fonttitle=\bfseries, arc=2mm, left=3mm, right=3mm, top=2mm, bottom=2mm]

\textbf{Scenario 2 --- Bottling Plant Quality Control}\\
A factory bottles fruit juice. The label states 12 fl oz per bottle. Quality control
tests $H_0{:}\ \mu = 12$ fl oz against $H_a{:}\ \mu \ne 12$ fl oz using $\alpha = 0.05$.

\begin{enumerate}[label=\alph*.]
  \item State the hypotheses in symbols. \writeline

  \item Describe a Type I error in context. \writelines{2}

  \item Describe a Type II error in context. \writelines{2}

  \item Which error is more costly for the company? Which is more costly for
        consumers? Explain. \writelines{3}
\end{enumerate}
\end{tcolorbox}

\begin{tcolorbox}[colback=white, colframe=black!40, title=\textbf{Tier E --- Extension},
    fonttitle=\bfseries, arc=2mm, left=3mm, right=3mm, top=2mm, bottom=2mm]

\textbf{Scenario 3 --- Clinical Drug Trial}\\
A pharmaceutical company tests whether a new antibiotic reduces mean infection duration
below 7 days. They set $\alpha = 0.01$ to minimize the risk of approving an ineffective drug.

\begin{enumerate}[label=\alph*.]
  \item Describe both error types in context and identify the consequences of each. \writelines{4}

  \item The FDA reviewer argues that $\alpha$ should be raised to 0.10 to avoid missing
        effective drugs. Write a counterargument explaining the trade-off. \writelines{3}

  \item If the sample size is increased from $n = 50$ to $n = 200$ (with $\alpha = 0.01$
        unchanged), what happens to the Type II error rate and power? Explain. \writelines{3}
\end{enumerate}
\end{tcolorbox}

\end{document}
